A transmissão de sinais em banda dupla concorrente tem se destacado por suprir as altas demandas de taxas de dados, proporcionando eficiência espectral. Contudo, a elevada relação de potência de pico por potência média (PAPR) característica desses sinais compromete o rendimento dos amplificadores de potência (PA). A aplicação de uma saturação digital controlada (ceifamento) surge como alternativa para operar o PA em regiões de maior eficiência, embora introduza distorções não lineares e espalhamento espectral que tendem a violar as máscaras normativas. Este trabalho propõe a aplicação de uma técnica de saturação para sinais de banda dupla (padrões LTE e Wi-Fi) associada a uma filtragem digital otimizada. Para assegurar a viabilidade de implementação em \textit{hardware}, o projeto focou na minimização da complexidade computacional por meio da redução da ordem do filtro. Utilizou-se o método de Powell para otimizar iterativamente os coeficientes de um filtro FIR com comprimento restrito a 9 amostras ($C=9$), partindo de um vetor inicial baseado na resposta ao impulso ideal das máscaras das normas. Os resultados demonstraram que o filtro otimizado conteve o espalhamento espectral dentro dos limites regulamentares, preservando a integridade de fase e amplitude do sinal original. A abordagem manteve a \textit{Error Vector Magnitude} (EVM) em níveis aceitáveis e consolidou uma redução da PAPR.
	